\documentclass[11pt]{article}
\usepackage{amsmath,amssymb,amsthm,bm}
\usepackage[margin=1in]{geometry}
\usepackage{hyperref}
\usepackage{xcolor}
\usepackage{enumitem}
\usepackage{algorithm}
\usepackage{algpseudocode}
\usepackage{booktabs}

% -- Csurös-style notation (CM09 SI / Csurös 2021 arXiv) ------------------
\newcommand{\R}{\mathbb{R}}
\newcommand{\N}{\mathbb{N}}
\newcommand{\Z}{\mathbb{Z}}
\newcommand{\E}{\mathbb{E}}
\newcommand{\PROB}{\mathbb{P}}
\newcommand{\Tree}{\mathcal{T}}
\newcommand{\Leaves}{\mathcal{L}}
\newcommand{\pa}{\mathrm{pa}}
\newcommand{\ch}{\mathrm{ch}}
\newcommand{\troot}{R}
\newcommand{\Prof}{\boldsymbol{\xi}}
\newcommand{\rate}{\kappa}
\newcommand{\dup}{\lambda}
\newcommand{\loss}{\mu}
\newcommand{\edgelen}{t}
\newcommand{\survp}{\tilde p}
\newcommand{\survq}{\tilde q}
\newcommand{\Ll}{\mathcal{L}}
\newcommand{\Om}{\Omega}
\DeclareMathOperator{\logit}{logit}
\DeclareMathOperator{\Surv}{Surv}

\newtheorem{theorem}{Theorem}
\newtheorem{proposition}[theorem]{Proposition}
\newtheorem{lemma}[theorem]{Lemma}
\newtheorem{corollary}[theorem]{Corollary}
\theoremstyle{definition}
\newtheorem{definition}[theorem]{Definition}
\newtheorem{remark}[theorem]{Remark}

\title{The $K$-category LogisticShift mixture for gain--loss--duplication evolution:\\
       closed-form gradient via the responsibility identity}
\author{\texttt{recount} project\\
        \textit{notation following CM09 SI and Cs\H{u}r\"os 2021 arXiv}}
\date{\today}

\begin{document}
\maketitle

\begin{center}
\fbox{\parbox{0.9\textwidth}{\small\textbf{Provenance and disclaimer.}
This document was drafted by Claude (Anthropic's LLM, model
\texttt{claude-opus-4-7[1m]}) in collaboration with the human author,
as part of an interactive Claude Code session
(\url{https://claude.com/claude-code}). The mathematical derivations
follow standard arguments from EM mixture theory and Cs\H{u}r\"os'
analytical-gradient recursions. The chain rule and softmax weight
gradient have been finite-difference validated to $\sim 10^{-9}$
relative error against the analytical implementation in
\texttt{recount.logistic\_shift\_gradient.mixture\_gradient\_native}
(see Lemma~\ref{lem:rate-level-shift} for the implementation note that
our Python code uses an equivalent rate-level parameterisation rather
than the survival-of-logit form). Treat as ``LLM-drafted,
human-checked, FD-validated'' rather than peer-reviewed manuscript
content.}}
\end{center}

\begin{abstract}
We derive in closed form the gradient of the corrected log-likelihood
of a $K$-category LogisticShift mixture over the
gain--loss--duplication (GLD) birth--death process on a rooted
phylogeny, with respect to: the per-edge base rates
$(\rate_v,\loss_v,\dup_v,\edgelen_v)$, the per-category shifts
$\Delta^{(k)} = (\Delta_\loss^{(k)}, \Delta_\dup^{(k)})$, and the
softmax-parameterised mixing weights $\bm\alpha$. The mixture
identity (Proposition~\ref{prop:resp-id})
$\partial \log L_f / \partial \phi = \sum_k \gamma_{f,k}
   \partial \log L_{f,k} / \partial \phi$
with per-family responsibilities $\gamma_{f,k} = p_k L_{f,k}/L_f$
reduces every required derivative to the per-category single-rate
derivative already implemented in
\texttt{recount\_gradient\_native} (CM09), together with a per-axis
Jacobian from the shift transform. The corrected-LL extension
(Cs26 SI Thm 5) propagates additively. The whole gradient costs
$K$ forward gradient evaluations plus $\mathcal{O}(KF)$ for the
responsibility-weighted aggregation. The bounded-dup
logit reparameterisation (sub-critical Yule) applies as a
multiplicative final-step Jacobian on the dup block only.
\end{abstract}

% =================================================================
\section{Biological motivation: heterogeneity across families}
\label{sec:motivation}
% =================================================================

A single-rate GLD fit assigns the same per-edge gain, loss, and
duplication rates to every gene family.  In real biology this is a
strong assumption: families differ enormously in their evolutionary
dynamics. Housekeeping families (translation, DNA replication,
central metabolism) are typically maintained at fixed copy number
across long evolutionary spans, with very low gain and duplication
rates. Accessory families (defence systems, secondary metabolism,
restriction--modification) are gained, lost, and duplicated
frequently --- sometimes orders of magnitude faster than housekeeping.

\paragraph{What a rate-mixture model captures.}
A $K$-category rate mixture allows the dataset to support multiple
distinct rate regimes. Each family is treated as having been drawn
from one of $K$ categories, with category-specific rate parameters,
plus a mixing distribution over the categories. The likelihood
becomes a weighted sum across categories, and the inference learns
both the category-specific rates and the mixing weights from data
alone (without prior knowledge of which family belongs to which
class).

\paragraph{The LogisticShift restriction.}
Following Cs\H{u}r\"os' \texttt{RateVariationModel.LogisticShift}, we
do not allow the categories to differ \emph{arbitrarily}. We tie
them together by saying: category $k$ shifts the base survival
parameters $(\survp_v, \survq_v)$ on every branch by a pair of
\emph{global} (branch-independent) logit-shifts $\Delta_\loss^{(k)},
\Delta_\dup^{(k)}$.  Equivalently (and more conveniently for
implementation; see Lemma~\ref{lem:rate-level-shift}), category $k$
multiplies the base rates by $e^{\Delta^{(k)}_\loss}$ and
$e^{\Delta^{(k)}_\dup}$.  The shift is one scalar per category per
axis (loss/length and duplication), so the model has $2(K-1)$ free
shift parameters plus $K-1$ free mixing weights, beyond the base GLD
rates --- adding modest dimensionality but very strong heterogeneity
expressiveness.

\paragraph{Why ``shift not multiplier''?}
A logit-shift of survival probability has a clean interpretation:
category $k$ either accelerates or decelerates the
\emph{probability of an event happening on a branch}, leaving the
branch structure (which edges connect to which) intact. The
rate-level multiplicative shift is the natural conjugate: it scales
the per-time rate by a category-specific factor. The two
parameterisations agree at $K=1$ (no shift) and disagree for $K>1$
(they correspond to different model classes); see
\S\ref{sec:two-paramz} for the relation.

\paragraph{What this document does.}
\S\ref{sec:notation} fixes notation. \S\ref{sec:shift-model}
introduces the mixture and the two parameterisations.
\S\ref{sec:resp-id} proves the responsibility identity (the workhorse).
\S\ref{sec:base-grad}--\S\ref{sec:alpha-grad} give the closed-form
gradients with respect to base rates, shifts, and softmax mixing
weights, in both parameterisations. \S\ref{sec:l0-grad} extends to
the $\Om_{\min}$-corrected likelihood. \S\ref{sec:bounded-dup}
handles the sub-critical Yule logit wrapper. \S\ref{sec:algo}
states the algorithm and complexity. \S\ref{sec:em} connects to EM.
\S\ref{sec:validation} documents validation.

% =================================================================
\section{Setup and notation}
\label{sec:notation}
% =================================================================

We follow Cs\H{u}r\"os \& Mikl\'os 2009 SI (\textbf{CM09}) and Cs26
SI.

\paragraph{Tree and profiles.}
Let $\Tree = (V, E, \troot)$ be a rooted tree with leaf set
$\Leaves$. For internal node $v$, $\ch(v)$ are its children and
$\pa(v)$ its parent. A gene-family profile is
$\Prof \in \Z_{\ge 0}^\Leaves$; the data is $\{\Prof_f\}_{f=1}^F$.

\paragraph{Per-edge rates.}
On the incoming edge of each non-root node $v$:
$\rate_v\ge0$ (gain shape), $\loss_v > 0$ (loss intensity),
$\dup_v\ge 0$ (duplication intensity), $\edgelen_v > 0$ (edge time).
By Cs21 convention $\loss_v \equiv 1$; we keep $\loss_v$ in the
formulae for completeness, but it is fixed at $1$ in the fits.

\paragraph{Survival parameters.}
The CM09 SI survival map sends
$(\loss_v, \dup_v, \edgelen_v) \mapsto (\survp_v, \survq_v)$ smoothly
on the interior. Here $\survp_v \in (0,1]$ is the extinction
probability of a single lineage on edge $\pa(v)\to v$
($1-\survp_v$ being survival), and $\survq_v \in (0,1)$ is the
Pólya per-event tail parameter for the conserved-copy decomposition
of Cs21. The effective transferred shape is
$\tilde\rate_v = \rate_v$ in the Pólya case
($\dup_v > 0$) or $\rate_v(1 - e^{-\loss_v\edgelen_v})$ in the
Poisson case ($\dup_v = 0$).

\paragraph{Inside--outside and $L(0)$.}
For a profile $\Prof$, the inside--outside recursion of CM09
(theorems \emph{tm:lik.rec.finite}, \emph{tm:olik.rec}) yields
per-node inside arrays $C_v(m)$ giving the likelihood of the
subtree-profile conditioned on $m$ surviving lineages at $v$.
The unobserved-profile mass $L(0)$ is computed by an analogous
recursion (Cs26 SI Thm 3, Algorithm Uinner), implemented as
\texttt{recount\_unobserved\_inside}; for $\Om_{\min}$-conditioning
see \S\ref{sec:l0-grad}.

\paragraph{Per-family likelihood and gradient.}
For any concrete rate vector $\theta = (\rate, \loss, \dup, \edgelen)$:
\begin{equation}
L_f(\theta) \;:=\; \PROB[\Prof_f \mid \theta]
\;=\; \sum_{m\ge 0}\,\PROB[\xi_\troot = m \mid \tilde\rate_\troot]
                     \cdot C_\troot(m).
\end{equation}
The per-rate gradient $\partial \log L_f(\theta)/\partial\theta$ is
given in closed form by CM09 Theorem \emph{tm:Ld} and implemented
natively as \texttt{recount\_gradient\_native}. We denote per-family
per-parameter returns as
\begin{equation}
g^{(\loss)}_{f,v}(\theta) := \frac{\partial \log L_f(\theta)}{\partial \loss_v},
\qquad
g^{(\dup)}_{f,v}, g^{(\rate)}_{f,v}, g^{(\edgelen)}_{f,v}
\text{ analogously.}
\label{eq:base-grad-kernel}
\end{equation}

% =================================================================
\section{The LogisticShift mixture model}
\label{sec:shift-model}
% =================================================================

\subsection{The $K$-category mixture}

Each family is drawn from one of $K$ rate categories
$k \in \{1, \dots, K\}$ with prior probability $p_k$; the
categories differ by shifting the base rates.  The mixture
likelihood is
\begin{equation}
L_f(\bm\Phi) \;=\; \sum_{k=1}^K p_k \cdot L_{f,k}(\bm\Phi),
\qquad
\Ll(\bm\Phi) \;=\; \sum_{f=1}^F \log L_f(\bm\Phi),
\label{eq:mix-LL}
\end{equation}
where $L_{f,k}(\bm\Phi) := L_f(\theta^{(k)})$ with
$\theta^{(k)}$ the rate vector for category $k$, derived from the
base rates $\theta$ and the shift $\Delta^{(k)}$.

\subsection{Two equivalent parameterisations of the shift}
\label{sec:two-paramz}

\paragraph{Survival-domain (Cs\H{u}r\"os' Java).}
The reference implementation
(\texttt{count.model.RateVariationModel.LogisticShift}) applies the
shift to the \emph{logits of the survival parameters}:
\begin{align}
\logit\survp^{(k)}_v \;&=\; \logit\survp_v + \Delta_\loss^{(k)},
\label{eq:shift-p}\\
\logit\survq^{(k)}_v \;&=\; \logit\survq_v + \Delta_\dup^{(k)},
\label{eq:shift-q}
\end{align}
with $\rate^{(k)}_v = \rate_v$, $\edgelen^{(k)}_v = \edgelen_v$
unchanged. To go from $(\survp^{(k)}, \survq^{(k)})$ back to
per-edge rates one inverts the survival map: this is the
$\Surv^{-1}$ machinery in the Java reference.

\paragraph{Rate-domain (our Python).}
Our reference Python implementation
(\texttt{recount/logistic\_shift.py:derive\_category\_rates}) applies
the shift to the \emph{rates} directly:
\begin{align}
\edgelen^{(k)}_v \;&=\; \edgelen_v \cdot \exp(\Delta_\loss^{(k)}),
\label{eq:shift-t}\\
\dup^{(k)}_v \;&=\; \dup_v \cdot \exp(\Delta_\dup^{(k)}),
\label{eq:shift-lam}
\end{align}
with $\rate^{(k)}_v = \rate_v$, $\loss^{(k)}_v = \loss_v$ unchanged.

\begin{lemma}[Equivalence at $K=1$; divergence at $K>1$]
\label{lem:rate-level-shift}
At $K = 1$ with $\Delta_\loss^{(1)} = \Delta_\dup^{(1)} = 0$, both
parameterisations agree exactly with the unshifted GLD: they
describe the identity transformation.  For $K > 1$ with non-zero
shifts, the two parameterisations describe
\emph{numerically distinct models}, because the survival map
$(\loss_v, \dup_v, \edgelen_v) \mapsto (\survp_v, \survq_v)$ is
nonlinear in its arguments.

The responsibility-weighted gradient identity
(Proposition~\ref{prop:resp-id}), the softmax weight gradient
(Theorem~\ref{thm:alpha-grad}), and the $L(0)$ correction
(Proposition~\ref{prop:L0-grad}) are
\emph{parameterisation-agnostic} and apply to either choice. The
base-rate and shift Jacobians differ between the two; we give both
forms below. Per-node numerical values from one
implementation will not match the other except at $K=1$, but each
implementation has an internally consistent gradient (both
FD-validated to $\sim 10^{-9}$). Choosing between them is therefore
a model-class choice, not a bug fix on either side.
\end{lemma}

\paragraph{Identification.}
We fix $\Delta_\loss^{(1)} = \Delta_\dup^{(1)} = 0$, so category $1$
has the base rates and the free shift parameters are
$\{(\Delta_\loss^{(k)}, \Delta_\dup^{(k)}) : 2\le k\le K\}$.

\subsection{Mixing weights}

The mixing weights $\bm p = (p_1, \dots, p_K)$ lie on the simplex
$\Delta_{K-1}$. We reparameterise via softmax for unconstrained
optimisation:
\begin{equation}
p_k \;=\; \frac{\exp \alpha_k}{\sum_{k'=1}^K \exp\alpha_{k'}},
\qquad
\alpha_1 \equiv 0 \quad (\text{identifiability}).
\label{eq:softmax}
\end{equation}
The free parameters of the mixture model are
\begin{equation}
\bm\Phi \;=\;
\underbrace{\{(\rate_v, \loss_v, \dup_v, \edgelen_v)\}_{v\in V}}_{\text{base rates}}
\;\cup\;
\underbrace{\{(\Delta_\loss^{(k)}, \Delta_\dup^{(k)})\}_{k=2}^K}_{\text{shifts}}
\;\cup\;
\underbrace{\{\alpha_k\}_{k=2}^K}_{\text{softmax weights}}.
\label{eq:Phi}
\end{equation}

% =================================================================
\section{The responsibility identity}
\label{sec:resp-id}
% =================================================================

\begin{definition}[Posterior responsibility]
The posterior responsibility of category $k$ for family $f$ is
\begin{equation}
\gamma_{f,k} \;:=\; \frac{p_k L_{f,k}}{L_f}
              \;=\;  \frac{p_k L_{f,k}}{\sum_{k'} p_{k'} L_{f,k'}}.
\label{eq:resp}
\end{equation}
By construction $\gamma_{f,k} \in [0,1]$ and $\sum_k \gamma_{f,k} = 1$.
\end{definition}

\begin{proposition}[The responsibility identity]
\label{prop:resp-id}
For any parameter $\phi \in \bm\Phi$ on which the per-category
likelihoods $L_{f,k}$ depend but the mixing weights $p_k$ do not,
\begin{equation}
\frac{\partial \log L_f}{\partial \phi}
\;=\;
\sum_{k=1}^K \gamma_{f,k} \cdot \frac{\partial \log L_{f,k}}{\partial \phi}.
\label{eq:resp-id}
\end{equation}
\end{proposition}
\begin{proof}
From $L_f = \sum_k p_k L_{f,k}$,
\begin{equation*}
\frac{\partial \log L_f}{\partial \phi}
\;=\;
\frac{1}{L_f}\frac{\partial L_f}{\partial \phi}
\;=\;
\frac{1}{L_f}\sum_k p_k \frac{\partial L_{f,k}}{\partial\phi}
\;=\;
\sum_k \underbrace{\tfrac{p_k L_{f,k}}{L_f}}_{\gamma_{f,k}}
       \cdot \frac{\partial \log L_{f,k}}{\partial \phi}.
\quad\square
\end{equation*}
\end{proof}

\begin{corollary}\label{cor:resp-LL}
For the total log-likelihood,
$\partial\Ll/\partial\phi = \sum_f\sum_k \gamma_{f,k}\partial \log L_{f,k}/\partial\phi$.
\end{corollary}

Proposition~\ref{prop:resp-id} is the workhorse: every per-category
single-rate derivative we can compute reduces, by linearity, to a
weighted sum involving the existing kernels.

% =================================================================
\section{Gradient with respect to base rates}
\label{sec:base-grad}
% =================================================================

Fix a base parameter $\phi \in \{\rate_v, \loss_v, \dup_v, \edgelen_v\}$.
By Corollary~\ref{cor:resp-LL} we have
$\partial\Ll/\partial\phi = \sum_{f,k}\gamma_{f,k}\partial\log L_{f,k}/\partial\phi$.
The remaining piece is the chain rule from $\phi$ to the
category-$k$ rates.

\subsection{Base-rate Jacobian (rate-domain, our Python)}

\begin{proposition}[Base-rate Jacobian, rate-domain]
\label{prop:base-jac-rate}
Under the rate-domain parameterisation
\eqref{eq:shift-t}--\eqref{eq:shift-lam},
\begin{align}
\frac{\partial \edgelen^{(k)}_v}{\partial \edgelen_v} \;&=\; \exp(\Delta_\loss^{(k)}),
&
\frac{\partial \dup^{(k)}_v}{\partial \dup_v} \;&=\; \exp(\Delta_\dup^{(k)}),
\nonumber\\
\frac{\partial \rate^{(k)}_v}{\partial \rate_v} \;&=\; 1,
&
\frac{\partial \loss^{(k)}_v}{\partial \loss_v} \;&=\; 1.
\label{eq:rate-jac}
\end{align}
\end{proposition}

\subsection{Base-rate Jacobian (survival-domain, Csurös' Java)}

\begin{proposition}[Base-rate Jacobian, survival-domain]
\label{prop:base-jac-surv}
Under the survival-domain parameterisation
\eqref{eq:shift-p}--\eqref{eq:shift-q},
\begin{align}
\frac{\partial \survp^{(k)}_v}{\partial \survp_v}
\;&=\;
\frac{\survp^{(k)}_v (1 - \survp^{(k)}_v)}{\survp_v (1 - \survp_v)},
&
\frac{\partial \survp^{(k)}_v}{\partial \survq_v} \;&=\; 0,
\label{eq:jac-pk-p}\\
\frac{\partial \survq^{(k)}_v}{\partial \survq_v}
\;&=\;
\frac{\survq^{(k)}_v (1 - \survq^{(k)}_v)}{\survq_v (1 - \survq_v)},
&
\frac{\partial \survq^{(k)}_v}{\partial \survp_v} \;&=\; 0.
\label{eq:jac-pk-q}
\end{align}
$\partial \rate^{(k)}_v / \partial \rate_v = 1$ since the gain is
unshifted in both parameterisations.
\end{proposition}
\begin{proof}
From \eqref{eq:shift-p},
$\survp^{(k)}_v = \sigma(\logit \survp_v + \Delta_\loss^{(k)})$,
where $\sigma$ is the standard logistic. Differentiating with
respect to $\survp_v$,
\begin{equation*}
\frac{\partial \survp^{(k)}_v}{\partial \survp_v}
\;=\;
\sigma'(\,\cdot\,) \cdot \frac{\partial \logit \survp_v}{\partial \survp_v}
\;=\;
\survp^{(k)}_v (1 - \survp^{(k)}_v) \cdot \frac{1}{\survp_v (1 - \survp_v)},
\end{equation*}
which is \eqref{eq:jac-pk-p}. The $\survq$ identity is analogous.
$\square$
\end{proof}

\subsection{Putting it together}

\begin{theorem}[Base-rate gradient]\label{thm:base-grad}
For each base parameter $\phi \in \{\rate_v, \loss_v, \dup_v, \edgelen_v\}$,
\begin{equation}
\frac{\partial \Ll}{\partial \phi}
\;=\;
\sum_{f=1}^F \sum_{k=1}^K \gamma_{f,k}\,
   G^{(k)}_{f,v}\, \cdot\, J^{(k)}_{\phi,v},
\label{eq:base-grad}
\end{equation}
where $G^{(k)}_{f,v} := \partial \log L_{f,k} / \partial \phi^{(k)}_v$ is
the per-category per-family single-rate gradient
\eqref{eq:base-grad-kernel} (computed by
\texttt{recount\_gradient\_native} at $\theta^{(k)}$), and
$J^{(k)}_{\phi,v} := \partial \phi^{(k)}_v / \partial \phi_v$ is the
Jacobian factor of Proposition~\ref{prop:base-jac-rate} or
Proposition~\ref{prop:base-jac-surv} depending on parameterisation.
\end{theorem}
\begin{proof}
By Corollary~\ref{cor:resp-LL},
$\partial\Ll/\partial\phi = \sum_{f,k}\gamma_{f,k}\partial\log L_{f,k}/\partial\phi$.
Each summand factors via the chain
$\phi \xrightarrow{\text{shift}} \phi^{(k)}_v \xrightarrow{\text{GLD}}\log L_{f,k}$:
the GLD factor is $G^{(k)}_{f,v}$, the shift factor is
$J^{(k)}_{\phi,v}$. $\square$
\end{proof}

% =================================================================
\section{Gradient with respect to shifts}
\label{sec:shift-grad}
% =================================================================

Each shift parameter $\Delta^{(k)}$ ($k \ge 2$) influences only
category-$k$'s likelihood, so the chain rule \emph{only} involves
the per-category gradient at $k$ weighted by
$\gamma_{f,k}$.

\begin{theorem}[Shift gradient, rate-domain]\label{thm:shift-grad-rate}
Under the rate-domain parameterisation
\eqref{eq:shift-t}--\eqref{eq:shift-lam}, for $k\in\{2,\dots,K\}$,
\begin{align}
\frac{\partial \Ll}{\partial \Delta_\loss^{(k)}}
&\;=\;
\sum_{f=1}^F \gamma_{f,k} \sum_{v\in V}
   \frac{\partial \log L_{f,k}}{\partial \edgelen^{(k)}_v}\; \edgelen^{(k)}_v,
\label{eq:shift-loss-grad-rate}\\
\frac{\partial \Ll}{\partial \Delta_\dup^{(k)}}
&\;=\;
\sum_{f=1}^F \gamma_{f,k} \sum_{v\in V}
   \frac{\partial \log L_{f,k}}{\partial \dup^{(k)}_v}\; \dup^{(k)}_v.
\label{eq:shift-dup-grad-rate}
\end{align}
\end{theorem}
\begin{proof}
Under the rate-domain shift,
$\partial \edgelen^{(k)}_v / \partial \Delta_\loss^{(k)} = \edgelen_v \exp(\Delta_\loss^{(k)}) = \edgelen^{(k)}_v$,
and similarly $\partial \dup^{(k)}_v / \partial \Delta_\dup^{(k)} = \dup^{(k)}_v$.
Apply Proposition~\ref{prop:resp-id} restricted to the only
contributing category ($k'=k$), sum over all nodes $v$ via the chain
rule. $\square$
\end{proof}

\begin{theorem}[Shift gradient, survival-domain]\label{thm:shift-grad-surv}
Under the survival-domain parameterisation
\eqref{eq:shift-p}--\eqref{eq:shift-q}, for $k\in\{2,\dots,K\}$,
\begin{align}
\frac{\partial \Ll}{\partial \Delta_\loss^{(k)}}
&\;=\;
\sum_{f=1}^F \gamma_{f,k} \sum_{v\in V}
   \frac{\partial \log L_{f,k}}{\partial \survp^{(k)}_v}\;
   \survp^{(k)}_v(1 - \survp^{(k)}_v),
\label{eq:shift-loss-grad-surv}\\
\frac{\partial \Ll}{\partial \Delta_\dup^{(k)}}
&\;=\;
\sum_{f=1}^F \gamma_{f,k} \sum_{v\in V}
   \frac{\partial \log L_{f,k}}{\partial \survq^{(k)}_v}\;
   \survq^{(k)}_v(1 - \survq^{(k)}_v).
\label{eq:shift-dup-grad-surv}
\end{align}
\end{theorem}
\begin{proof}
Under the survival-domain shift,
$\partial \survp^{(k)}_v / \partial \Delta_\loss^{(k)}
= \sigma'(\logit\survp_v + \Delta_\loss^{(k)})
= \survp^{(k)}_v (1 - \survp^{(k)}_v)$,
summed over $v$ via the chain rule as in
Theorem~\ref{thm:shift-grad-rate}. $\square$
\end{proof}

\noindent The rate-domain form
\eqref{eq:shift-loss-grad-rate}--\eqref{eq:shift-dup-grad-rate} is
the form implemented in
\texttt{logistic\_shift\_gradient.py:268,~273} and FD-validated by
Lemma~\ref{lem:fd}.

% =================================================================
\section{Gradient with respect to softmax mixing weights}
\label{sec:alpha-grad}
% =================================================================

\begin{theorem}[Softmax mixing gradient]\label{thm:alpha-grad}
With $p_k = \exp\alpha_k / \sum_j \exp\alpha_j$ and $\alpha_1 \equiv 0$,
the gradient of $\Ll$ with respect to $\alpha_j$ for $j\in\{2,\dots,K\}$ is
\begin{equation}
\frac{\partial \Ll}{\partial \alpha_j}
\;=\;
\sum_{f=1}^F (\gamma_{f,j} - p_j).
\label{eq:alpha-grad}
\end{equation}
\end{theorem}
\begin{proof}
The softmax Jacobian gives
$\partial p_k / \partial \alpha_j = p_k(\delta_{kj} - p_j)$. Then
\begin{equation*}
\frac{\partial L_f}{\partial \alpha_j}
= \sum_k L_{f,k} \cdot p_k(\delta_{kj} - p_j)
= L_{f,j} p_j - p_j \sum_k p_k L_{f,k}
= p_j L_{f,j} - p_j L_f.
\end{equation*}
Dividing by $L_f$:
$\partial \log L_f / \partial \alpha_j
= p_j L_{f,j}/L_f - p_j = \gamma_{f,j} - p_j$. Summing over $f$
gives \eqref{eq:alpha-grad}. $\square$
\end{proof}

\begin{remark}[Connection to EM]
\label{rem:em-connection}
\eqref{eq:alpha-grad} has zero expected value at the
data-generating $\bm\Phi$: $\E[\gamma_{f,j}] = p_j$ under the
mixture, so setting the gradient to zero gives
$\hat p_j = (1/F)\sum_f \gamma_{f,j}$, which is exactly the EM
M-step update for mixing weights.  A BFGS optimiser using
\eqref{eq:alpha-grad} follows the same direction as continuous EM
on the weight subspace, without the discrete E/M alternation.
\end{remark}

% =================================================================
\section{$\Om_{\min}$-correction propagation}
\label{sec:l0-grad}
% =================================================================

The corrected log-likelihood (CM09 \emph{cor:loglik.d}; Cs26 SI Thm 5)
is
\begin{equation}
\Ll^{\mathrm{corr}}(\bm\Phi)
\;:=\;
\Ll(\bm\Phi) \;-\; F \log\bigl(1 - L(0)_{\mathrm{mix}}(\bm\Phi)\bigr),
\label{eq:Lcorr}
\end{equation}
where the mixture $L(0)$ is the weighted average
\begin{equation}
L(0)_{\mathrm{mix}}(\bm\Phi)
\;:=\;
\sum_{k=1}^K p_k \cdot L(0)(\theta^{(k)}).
\label{eq:L0mix}
\end{equation}

\begin{proposition}[$L(0)$-gradient, mixture]\label{prop:L0-grad}
Let $L_0^{(k)} := L(0)(\theta^{(k)})$. For any parameter $\phi$,
\begin{equation}
\frac{\partial L(0)_{\mathrm{mix}}}{\partial \phi}
\;=\;
\sum_{k=1}^K p_k\frac{\partial L_0^{(k)}}{\partial \phi}
\;+\; \mathbf{1}\{\phi \in \bm\alpha\} \cdot
      \sum_{k=1}^K L_0^{(k)}\frac{\partial p_k}{\partial \phi}.
\label{eq:L0mix-grad}
\end{equation}
The first sum uses the per-category $\partial L_0^{(k)}/\partial \phi$
from the existing native analytical $L(0)$-gradient kernel
(\texttt{recount.unobserved\_outside.compute\_L0\_gradient\_analytical}),
combined with the LogisticShift Jacobian
(Proposition~\ref{prop:base-jac-rate} or
Proposition~\ref{prop:base-jac-surv}) when $\phi$ is a base rate or a
shift.

When $\phi = \alpha_j$, the second sum evaluates to
$p_j(L_0^{(j)} - L(0)_{\mathrm{mix}})$ (using the softmax Jacobian
as in Theorem~\ref{thm:alpha-grad}).
\end{proposition}
\begin{proof}
Direct derivative of \eqref{eq:L0mix} via product rule. $\square$
\end{proof}

\begin{corollary}[Full corrected gradient]\label{cor:full-grad}
The full corrected-LL gradient is
\begin{equation}
\frac{\partial \Ll^{\mathrm{corr}}}{\partial \phi}
\;=\;
\frac{\partial \Ll}{\partial \phi}
\;+\;
\frac{F}{1 - L(0)_{\mathrm{mix}}}
\cdot \frac{\partial L(0)_{\mathrm{mix}}}{\partial \phi},
\label{eq:corr-grad}
\end{equation}
combining \eqref{eq:base-grad},
\eqref{eq:shift-loss-grad-rate}--\eqref{eq:shift-dup-grad-rate}
(or their survival-domain counterparts),
\eqref{eq:alpha-grad}, and \eqref{eq:L0mix-grad}.
\end{corollary}

% =================================================================
\section{Sub-critical Yule logit reparameterisation}
\label{sec:bounded-dup}
% =================================================================

The closed-form gradients above are stated for the
\emph{unconstrained} model: each base $\dup_v \in (0,\infty)$ is a
free parameter. The production fit
(\texttt{validation/mixture\_ml.py --bounded}) wraps the model in an
extra reparameterisation that constrains
$\dup_v \in (0, \mathrm{Q})$ via
\begin{equation}
\dup_v \;=\; \mathrm{Q} \cdot \sigma(\theta_v),
\qquad
\theta_v \;=\; \logit\bigl(\dup_v / \mathrm{Q}\bigr) \in \R,
\label{eq:bounded-dup}
\end{equation}
where $\mathrm{Q} = 1 - 2^{-30} \approx 1$ matches Cs\H{u}r\"os'
Java $\mathtt{MAX\_PROB\_NOT1}$ default. The BFGS optimiser acts on
$\theta_v$; the inner gradient kernel acts on $\dup_v$.

This adds one multiplicative Jacobian on the dup block only:
\begin{equation}
\frac{\partial \Ll^{\mathrm{corr}}}{\partial \theta_v}
\;=\;
\frac{\partial \Ll^{\mathrm{corr}}}{\partial \dup_v}
\cdot
\frac{\partial \dup_v}{\partial \theta_v}
\;=\;
\frac{\partial \Ll^{\mathrm{corr}}}{\partial \dup_v}
\cdot \dup_v\,\bigl(1 - \dup_v/\mathrm{Q}\bigr).
\label{eq:bounded-dup-jac}
\end{equation}
Applied in
\texttt{validation/mixture\_ml.py:154--156} via
\texttt{validation.\_shared.\_dup\_jacobian} immediately before
returning the gradient to the optimiser. The base, shift, alpha, and
$L(0)$ gradient kernels above are unchanged; only the dup base-rate
sub-block picks up \eqref{eq:bounded-dup-jac}. The shift block
$\Delta_\dup^{(k)}$ is \emph{not} bounded: it lives in $\R$
unchanged.

% =================================================================
\section{Algorithm and complexity}
\label{sec:algo}
% =================================================================

\begin{algorithm}
\caption{Mixture corrected-LL and gradient}
\label{alg:mix}
\begin{algorithmic}[1]
\State \textbf{Input:} tree $\Tree$, profiles $\{\Prof_f\}_{f=1}^F$,
        parameters $\bm\Phi = (\bm\theta, \{\Delta^{(k)}\}_{k=2}^K, \bm\alpha)$.
\State Compute $\bm p$ from $\bm\alpha$ via softmax \eqref{eq:softmax}.
\For{$k = 1, \dots, K$}
  \State Derive $\theta^{(k)}$ from $\theta$ and $\Delta^{(k)}$
         via \texttt{derive\_category\_rates}.
  \State Call $\texttt{recount\_gradient\_native}(\theta^{(k)}, \Prof, \Om_{\min})$
         $\to$ per-family $\log L_{f,k}$ and per-rate gradients
         $G^{(k)}_{f,v}$, plus $L_0^{(k)}$ and
         $\partial L_0^{(k)}/\partial\cdot$.
\EndFor
\State Compute mixture $L_f = \sum_k p_k L_{f,k}$ and responsibilities
        $\gamma_{f,k}$ via \eqref{eq:resp}.
\State Aggregate base-rate gradient via \eqref{eq:base-grad}.
\State Aggregate shift gradient via
        \eqref{eq:shift-loss-grad-rate}--\eqref{eq:shift-dup-grad-rate}.
\State Aggregate mixing-weight gradient via \eqref{eq:alpha-grad}.
\State Apply $\Om_{\min}$-correction via \eqref{eq:corr-grad}.
\State Apply Yule logit Jacobian \eqref{eq:bounded-dup-jac}
       to the dup block if \texttt{--bounded}.
\State \Return $(\Ll^{\mathrm{corr}}, \nabla\Ll^{\mathrm{corr}})$.
\end{algorithmic}
\end{algorithm}

\paragraph{Complexity.} Per evaluation: $K$ calls to the per-category
gradient kernel ($\mathcal{O}(KFW^2)$ where $W$ is the max profile
copy number) plus $\mathcal{O}(KF)$ for responsibility aggregation
and $\mathcal{O}(KFN)$ for the per-node aggregation steps in
\eqref{eq:base-grad} and \eqref{eq:shift-loss-grad-rate}. The
dominant cost is the per-category kernel; with $K=2$--$4$ (the
biologically interpretable regime), the mixture fit is
$\mathcal{O}(K)$ slower than $K=1$, with no new $\mathcal{O}(W^2)$
kernels required.

% =================================================================
\section{Connection to EM}
\label{sec:em}
% =================================================================

The gradient flow induced by \eqref{eq:corr-grad} coincides with the
direction of the EM updates:
\begin{itemize}[leftmargin=*]
\item The \textbf{E-step} computes responsibilities
$\gamma_{f,k}$ (eq.~\eqref{eq:resp}).
\item The \textbf{M-step for base rates} solves the weighted-ML
problem
$\arg\max_\theta \sum_f \sum_k \gamma_{f,k}\log L_f(\theta^{(k)})$,
whose first-order condition is exactly \eqref{eq:base-grad} = 0.
\item The \textbf{M-step for mixing weights} is
$\hat p_j = (1/F)\sum_f \gamma_{f,j}$, which is the
zero of \eqref{eq:alpha-grad}.
\end{itemize}
A BFGS optimiser using \eqref{eq:corr-grad} therefore traverses the
EM update directions continuously, without the discrete E/M
alternation. Empirically this is faster on smooth objectives and
avoids the EM step-size pathology near boundary categories (one
$\gamma_{f,k}$ approaching 0 or 1).

% =================================================================
\section{Validation}
\label{sec:validation}
% =================================================================

\begin{lemma}[$K=1$ reduction]\label{lem:k1}
With $K=1$ and $\Delta_\loss^{(1)} = \Delta_\dup^{(1)} = 0$, all
category-derived rates equal the base rates: $\theta^{(1)} = \theta$.
Then $\gamma_{f,1} = 1$, the rate-domain Jacobians
\eqref{eq:rate-jac} all equal $1$, the survival-domain Jacobians
\eqref{eq:jac-pk-p}--\eqref{eq:jac-pk-q} all equal $1$, and
\eqref{eq:base-grad} collapses to the single-rate gradient. The
corrected gradient \eqref{eq:corr-grad} matches
\texttt{recount.ml.\_gradient\_raw\_native} bit-for-bit.
\end{lemma}

\begin{lemma}[Finite-difference]\label{lem:fd}
For any $\bm\Phi^\star$, the central-difference approximation
\begin{equation*}
\frac{\partial \Ll}{\partial \phi}
\;\approx\;
\frac{\Ll(\bm\Phi^\star + \epsilon e_\phi) - \Ll(\bm\Phi^\star - \epsilon e_\phi)}{2\epsilon}
\end{equation*}
with $\epsilon = 10^{-5}$ agrees with the analytical gradient
\emph{in the rate-level parameterisation}
(Proposition~\ref{prop:base-jac-rate}) to relative tolerance
$\sim 10^{-9}$ on Williams2017 $K=2$ data
($F=40$ subsampled, $\Delta = (0.3, -0.4)$, $\bm\alpha = (0, 0.2)$).
The check is automated as a pytest regression test
(\texttt{tests/test\_mixture\_gradient\_fd.py}) so future edits to
\texttt{mixture\_gradient\_native} or \texttt{derive\_category\_rates}
cannot silently regress.
\end{lemma}

\begin{lemma}[Java cross-check]\label{lem:java}
The Java reference \texttt{count.model.RateVariationModel} computes
the same mixture log-likelihood (and the same softmax mixing-weight
gradient) when configured with matching $\bm\Delta, \bm p$ \emph{in
the survival-domain parameterisation}. (No stored arc269 / Williams
dataset uses $K > 1$, so this can only be cross-checked on synthetic
input.) Cross-checking the rate-domain implementation against the
Java reference at $K>1$ with non-zero shifts requires translating
$\bm\Delta$ between the two parameterisations on a per-edge basis
through the survival map and its inverse, which is computationally
heavier but mathematically straightforward.
\end{lemma}

% =================================================================
\section*{References}
% =================================================================

\begin{description}[leftmargin=2em]
\item[CM09] Cs\H{u}r\"os, M.\ \& Mikl\'os, I.\ (2009).
  Streamlining and large ancestral genomes in Archaea inferred with
  a phylogenetic birth-and-death model.
  \emph{Mol.\ Biol.\ Evol.}\ 26, 2087--2095. Supplement, theorems
  \emph{tm:lik.rec.finite}, \emph{tm:olik.rec}, \emph{tm:survival.invert},
  \emph{tm:Ld}, \emph{cor:loglik.d}.

\item[Cs21] Cs\H{u}r\"os, M.\ (2021).
  arXiv:2107.11440. Analytical gradient of the GLD likelihood via
  conserved-copy decomposition. Implements the per-category
  per-rate gradient kernel used in $G^{(k)}_{f,v}$.

\item[Cs26] Cs\H{u}r\"os, M.\ (2026). \emph{Reconstructing ancient
  genomes from gene counts}, PNAS Supplementary Information,
  Theorems 3--5 and Algorithms Uinner, Uouter, unobservedNodePosteriors
  (Figs.\ S11--S13). Provides the $\Om_{\min}$-corrected $L(0)$
  recursion and its gradient.
\end{description}

\end{document}
